\chapter{Chapter 2: The LLM Stack and Core Patterns}

\section*{Context}
\textbf{Problem statement:} Design a production LLM stack and decision framework for choosing between prompt chaining, RAG, and agentic orchestration under strict latency, cost, and safety constraints.


\subsection*{Real-World Examples}
\begin{itemize}
  \item A SaaS company routes invoice classification to a small model, customer policy Q&A to RAG, and dispute-resolution workflows to an agentic path.
  \item An enterprise migration uses managed APIs for launch, then moves stable high-volume routes to self-hosted GPUs for unit-cost reduction.
\end{itemize}

\section*{Step 1: Understand the Problem and Establish Design Scope}
\subsection*{Simulated Dialogue}
\textbf{Interviewer:} Design the LLM stack for a multi-product AI platform.\\
\textbf{Candidate:} I will quantify traffic, latency budgets, and risk classes first, then define pattern routing and serving architecture.

\subsection*{Functional Requirements}
\begin{itemize}
  \item Support prompt chaining, RAG, and agentic flows under one platform.
  \item Route requests dynamically across model tiers and providers.
  \item Enforce policy checks and deterministic fallback behavior.
  \item Provide traceable outputs with quality and cost telemetry.
\end{itemize}

\subsection*{Non-Functional Requirements}
\begin{itemize}
  \item Latency: route-specific TTFT/TPOT and p95 targets.
  \item Scale: absorb burst traffic without cascading failures.
  \item Cost: per-route token budgets and monthly controls.
  \item Reliability/compliance: tenant isolation, auditability, policy enforcement.
\end{itemize}

\subsection*{Back-of-the-envelope Estimation}
\begin{itemize}
  \item Requests/day: 16M; pattern split 50\% prompt chain, 35\% RAG, 15\% agentic.
  \item Average QPS: $16,000,000/86,400 \approx 185$; peak QPS (12x): $\approx 2,220$.
  \item Total input/day: $\approx 29.44B$ tokens; total output/day: $\approx 3.888B$ tokens.
  \item Cost/day: $(input/10^6)\cdot p_{in} + (output/10^6)\cdot p_{out}$.
  \item Latency envelope: prompt chain p95 1.5--2.5s, RAG p95 3--5s, agentic p95 6--12s.
  \item VRAM planning: 7B class $\sim 16$--$24GB$, 13B--34B $\sim 40$--$80GB$, 70B class requires multi-GPU parallel serving.
\end{itemize}

\section*{Step 2: Propose High-Level Design and Get Buy-In}
\subsection*{Core Architecture Components}
\begin{itemize}
  \item API gateway + policy layer
  \item Intent/complexity router
  \item Prompt runtime, retrieval runtime, and agent runtime
  \item Model router and serving runtime with batching and KV cache
  \item Safety/runtime verification plus observability/eval loop
\end{itemize}

\subsection*{Pipeline A: Stack Preparation / Data and Runtime Setup}
\begin{enumerate}
  \item Register prompts, schemas, tools, and policy rules.
  \item Build retrieval artifacts and model catalog tiers.
  \item Configure route budgets and fallback chains.
  \item Publish config with versioned rollback handles.
\end{enumerate}

\subsection*{Pipeline B: Runtime Query Execution}
\begin{enumerate}
  \item Validate request and classify intent/complexity.
  \item Route into prompt, RAG, or agentic execution graph.
  \item Execute with bounded token/time budgets and streaming.
  \item Apply safety checks and emit traces/eval metrics.
\end{enumerate}

\textit{[Insert Architecture Diagram: Show gateway and policy layer feeding a router that branches into prompt, RAG, and agentic runtimes; all converge via model router + serving runtime with continuous batching/KV cache, with safety validators and observability loop.]}

\subsection*{TikZ Architecture Diagram}
\begin{center}
\begin{tikzpicture}[node distance=1.6cm, auto, >=stealth]
  \node[draw, rounded corners] (client) {UserClient};
  \node[draw, rounded corners, right=of client] (gateway) {APIGateway};
  \node[draw, rounded corners, right=of gateway] (router) {IntentRouter};
  \node[draw, rounded corners, above right=0.8cm and 1.5cm of router] (rag) {RAGRuntime};
  \node[draw, rounded corners, below right=0.8cm and 1.5cm of router] (agent) {AgentRuntime};
  \node[draw, rounded corners, right=3.2cm of router] (model) {ModelRouter};
  \node[draw, rounded corners, below=of model] (safe) {SafetyVerifier};
  \draw[->] (client) -- (gateway);
  \draw[->] (gateway) -- (router);
  \draw[->] (router) -- (rag);
  \draw[->] (router) -- (agent);
  \draw[->] (rag) -- (model);
  \draw[->] (agent) -- (model);
  \draw[->] (model) -- (safe);
  \draw[->] (safe) |- (client);
\end{tikzpicture}
\end{center}

\section*{Step 3: Design Deep Dive}
\begin{enumerate}
  \item \textbf{Pattern Selection: RAG vs Agent}
  \item \textbf{Prompt Chaining Design}
  \item \textbf{Model Routing and Tiering}
\end{enumerate}

\subsection*{Trade-off Comparisons}
\begin{tabular}{|p{0.18\linewidth}|p{0.24\linewidth}|p{0.24\linewidth}|p{0.24\linewidth}|}
\hline
\textbf{Decision} & \textbf{Option A} & \textbf{Option B} & \textbf{Recommended} \\
\hline
Execution pattern & Prompt chaining only & RAG/agentic everywhere & Route by complexity and risk \\
\hline
Model hosting & Managed-only & Self-hosted-only & Hybrid managed + self-hosted \\
\hline
Routing policy & Static mapping & Dynamic rule-based & Dynamic with eval-gated updates \\
\hline
Caching & No cache & Single cache layer & Semantic + KV cache hierarchy \\
\hline
\end{tabular}

\subsection*{Production Bottlenecks and Mitigations}
\begin{itemize}
  \item Provider rate limits and retry storms: use token buckets, circuit breakers, and jittered retries.
  \item Router drift: track confusion matrix, false escalations/downgrades, and rollback quickly.
  \item Metadata DB pool exhaustion: read replicas + TTL cache + connection caps.
  \item Stale embeddings/index skew: embedding versioning, shadow tests, alias-based rollback.
  \item Agentic hallucination loops: step ceilings, loop-signature detection, forced terminal action.
  \item GPU OOM: strict token caps, dynamic batch shrink, and route downgrades under memory pressure.
\end{itemize}

\subsection*{Advanced Techniques}
\begin{itemize}
  \item Continuous batching
  \item KV caching
  \item Semantic caching with policy-aware invalidation
  \item Cross-encoder reranking
  \item Parent-child chunking
  \item LoRA adapters
\end{itemize}

\section*{Step 4: Wrap Up}
\textbf{Final design:} A policy-driven execution platform that routes each request to the cheapest safe path meeting quality and latency targets.

\textbf{Wrap-around concerns:}
\begin{itemize}
  \item Routing drift and provider concentration risk
  \item Prompt/runtime governance and release gates
  \item Operational debt from uncontrolled chain/tool growth
\end{itemize}

\section*{Step 5: Interviewer Curveballs}
\subsection*{Interviewer Follow-ups \& Frequently Asked Questions}
\textbf{Q1:} If only two patterns are allowed in v1, which and why?\\
\textbf{The Ideal Answer:} Prompt chaining + RAG for best latency-cost-quality coverage.\\
\textbf{Deep Dive \& Trade-offs:} Add agentic only after measured unresolved complexity.\\
\textbf{The Interviewer's Rubric:} Tests prioritization and scope discipline.

\textbf{Q2:} Trigger to migrate a route from managed to self-hosted?\\
\textbf{The Ideal Answer:} Sustained TCO and SLO advantage for multiple weeks plus ops readiness.\\
\textbf{Deep Dive \& Trade-offs:} Include infra and on-call costs, keep managed failover path.\\
\textbf{The Interviewer's Rubric:} Tests practical migration judgment.

\textbf{Q3:} Detect router harming quality?\\
\textbf{The Ideal Answer:} Monitor route confusion, quality pass-rate by tier, and false downgrade/upgrade rates.\\
\textbf{Deep Dive \& Trade-offs:} Use offline oracle labels + online A/B evaluation.\\
\textbf{The Interviewer's Rubric:} Tests routing as first-class system component.

\textbf{Q4:} Hidden bottleneck in this stack?\\
\textbf{The Ideal Answer:} Metadata/policy lookup path often dominates p95 under load.\\
\textbf{Deep Dive \& Trade-offs:} Use pooled connections, replicas, and short TTL caches.\\
\textbf{The Interviewer's Rubric:} Tests full-path bottleneck diagnosis.

\textbf{Q5:} Keep agentic path from becoming cost sink?\\
\textbf{The Ideal Answer:} Enforce strict step/token budgets and deterministic terminal actions.\\
\textbf{Deep Dive \& Trade-offs:} Track cost per successful task and enforce planner budgets.\\
\textbf{The Interviewer's Rubric:} Tests control-plane cost governance.

\textbf{Q6:} When does semantic cache become dangerous?\\
\textbf{The Ideal Answer:} Loose similarity thresholds + long freshness windows on dynamic data.\\
\textbf{Deep Dive \& Trade-offs:} Use route-specific thresholds and source-version invalidation.\\
\textbf{The Interviewer's Rubric:} Tests safe cache policy design.

\textbf{Q7:} How to set reranker depth?\\
\textbf{The Ideal Answer:} Tune candidate size for recall gain vs latency overhead by query class.\\
\textbf{Deep Dive \& Trade-offs:} Usually 20--80 candidates depending on corpus and intent.\\
\textbf{The Interviewer's Rubric:} Tests practical retrieval optimization.

\textbf{Q8:} Minimum weekly KPI set?\\
\textbf{The Ideal Answer:} TTFT p95, TPOT p95, cost/query, route distribution, Recall@K, citation precision, safety violations.\\
\textbf{Deep Dive \& Trade-offs:} Evaluate by route/version to catch drift early.\\
\textbf{The Interviewer's Rubric:} Tests multidimensional operational maturity.

\section*{End of Chapter}
\begin{itemize}
  \item \textbf{Key takeaway:} Tie architecture decisions to measurable SLOs, cost, and safety outcomes.
  \item \textbf{What to practice:} Defend one default and one fallback design with explicit trade-offs.
  \item \textbf{Interview drill:} Explain regression detection, rollback triggers, and risk communication.
\end{itemize}
